\section{Detailed Explanation of Algorithm~1}
\label{app:algorithm-one-explanation}

This section follows Algorithm~1 and its proof in \cite[Section~3]{12WSNE}.

\subsection{Preprocessing}

\begin{enumerate}
  \item Compute a Nash equilibrium \((\mathbf{x}^*, \mathbf{y}^*)\) of the zero-sum game \((\mathcal{R}, -\mathcal{R})\); this can be solved in polynomial time via linear programming.

        By the properties of zero-sum games, \(\mathbf{x}^*\) is the row player's \textbf{maximin strategy}:
        \[
          \mathcal{R}(\mathbf{x}^*, \mathbf{y}^*) = \max_{\mathbf{x} \in \Delta^n} \min_{\mathbf{y} \in \Delta^n} \mathcal{R}(\mathbf{x}, \mathbf{y})
        \]
        Hence for every column strategy \(\mathbf{y}\), \(\mathcal{R}(\mathbf{x}^*, \mathbf{y}^*) \le \mathcal{R}(\mathbf{x}^*, \mathbf{y})\). \textbf{This property will be used several more times later.}

  \item Compute a Nash equilibrium \((\hat{\mathbf{x}}, \hat{\mathbf{y}})\) of the zero-sum game \((-\mathcal{C}, \mathcal{C})\). Similarly, for every row strategy \(\mathbf{x}\), \(\mathcal{C}(\hat{\mathbf{x}}, \hat{\mathbf{y}}) \le \mathcal{C}(\mathbf{x}, \hat{\mathbf{y}})\).
\end{enumerate}

We then case-analyze the relationship between \(\mathcal{R}(\mathbf{x}^*, \mathbf{y}^*)\) and \(\mathcal{C}(\hat{\mathbf{x}}, \hat{\mathbf{y}})\). By symmetry, we only analyze the case \textbf{\(\mathcal{R}(\mathbf{x}^*, \mathbf{y}^*) \ge \mathcal{C}(\hat{\mathbf{x}}, \hat{\mathbf{y}})\)}:

\subsection{The \texorpdfstring{\(\mathcal{R}(\mathbf{x}^*, \mathbf{y}^*) \ge \mathcal{C}(\hat{\mathbf{x}}, \hat{\mathbf{y}})\)}{R(x*,y*) >= C(xh,yh)} branch}

\paragraph{Case 1(a): \(\mathcal{R}(\mathbf{x}^*, \mathbf{y}^*) \le 1/2\)}
In this case, we simply output \((\hat{\mathbf{x}}, \mathbf{y}^*)\). Its correctness proof is given in Appendix~\ref{app:case-one-low}.

\paragraph{Case 1(b): ``one high, one low''}
In this case, \(\mathcal{R}(\mathbf{x}^*, \mathbf{y}^*) > 1/2\), and the following linear feasibility system has a solution \(\mathbf{x}'\):
\[
  \begin{cases}
    \sum_{i \in \supp(\mathbf{x}^*)} \mathbf{x}'(i) \cdot \mathcal{C}_{ij} \le \frac12, & \forall j \in [n]                 \\[4pt]
    \mathbf{x}'(i) \ge 0,                                                               & \forall i \in \supp(\mathbf{x}^*) \\[4pt]
    \sum_{i \in \supp(\mathbf{x}^*)} \mathbf{x}'(i) = 1
  \end{cases}
\]
This linear program can be solved, or declared infeasible, in polynomial time.
If it is feasible, we output \((\mathbf{x}', \mathbf{y}^*)\).

Its correctness proof is given in Appendix~\ref{app:case-one-low-high}.

\paragraph{Case 1(c): ``both high''}
In this case, neither of the previous two conditions holds; that is, we cannot bound a single player's best response by \(\le 1/2\).

Define the \textbf{restricted game} \((\mathcal{R}_{\mathbf{x}^*}, \mathcal{C}_{\mathbf{x}^*})\): the row player may only use actions in \(\supp(\mathbf{x}^*)\), while the column player may use all \(n\) actions.

One can show that every Nash equilibrium \((\mathbf{x}, \mathbf{y})\) of \((\mathcal{R}_{\mathbf{x}^*}, \mathcal{C}_{\mathbf{x}^*})\) satisfies:
\[
  \mathcal{R}(\mathbf{x}, \mathbf{y}) > \frac12,\;\; \mathcal{C}(\mathbf{x}, \mathbf{y}) > \frac12
\]
The contradiction argument is given in Appendix~\ref{app:case-one-both-high}.

Hence \((\mathbf{x}, \mathbf{y})\) is a \(1/2\)-WSNE of the original game. So we only need to compute such an \((\mathbf{x}, \mathbf{y})\). However, we cannot in general do this in polynomial time using known methods, because computing a NE of a bimatrix game is \PPAD-complete \cite{DGP09}.

Here we use a different sampling theorem that preserves the equilibrium payoffs on the sampled supports:

\begin{lemma}[{Constant-Support Well-Supporting Lemma \cite[Theorem~2]{cfj14}}]
  For any Nash equilibrium \((\mathbf{x}, \mathbf{y})\) of the restricted game \((\mathcal{R}_{\mathbf{x}^*}, \mathcal{C}_{\mathbf{x}^*})\), there exists a \(\kappa(\delta)\)-uniform strategy profile \((\mathbf{w}, \mathbf{z})\) such that:
  \[
    \forall i \in \supp(\mathbf{w}),\; \mathcal{R}(e_i, \mathbf{z}) \ge \mathcal{R}(\mathbf{x}, \mathbf{y}) - \delta > \frac12 - \delta
  \]
  \[
    \forall j \in \supp(\mathbf{z}),\; \mathcal{C}(\mathbf{w}, e_j) \ge \mathcal{C}(\mathbf{x}, \mathbf{y}) - \delta > \frac12 - \delta
  \]
  Since the maximum possible payoff is \(1\) (the payoff matrices are normalized), the above conditions imply:
  \[
    \forall i \in \supp(\mathbf{w}),\ \mathcal{R}(e_i, \mathbf{z}) \ge \left(\frac12 - \delta\right) - \left(1 - \max_k \mathcal{R}(e_k, \mathbf{z})\right) = \max_k \mathcal{R}(e_k, \mathbf{z}) - \left(\frac12 + \delta\right)
  \]
  and similarly for the column player. Therefore, \((\mathbf{w}, \mathbf{z})\) is a \((1/2+\delta)\)-WSNE.
\end{lemma}

Here \(\kappa(\delta) = \left\lceil \frac{2 \ln (1/\delta)}{\delta^2} \right\rceil\) is a constant for fixed \(\delta \in (0,1)\). A \(k\)-uniform strategy profile \((\mathbf{w}, \mathbf{z})\) satisfies that every component of \(\mathbf{w}, \mathbf{z}\) is a non-negative integer multiple of \(\frac{1}{k}\), so it naturally follows that \(\max(|\supp(\mathbf{w})|, |\supp(\mathbf{z})|) \le \kappa(\delta)\), i.e., both supports have size at most \(\kappa(\delta)\). There are \(O(n^{\kappa(\delta)})\) such strategies per player and hence \(O(n^{2\kappa(\delta)})\) profiles. Since the lemma guarantees existence, it suffices to enumerate all profiles and check whether each one is a \((1/2+\delta)\)-WSNE. When \(\delta\) is a constant, the running time of the algorithm is polynomial in \(n\).

\begin{remark}
  In Lemma~\ref{lem:lmm_thm}, the sampling size \(k = O\left(\frac{\log n}{\varepsilon^2}\right)\) depends on \(n\), while here \(\kappa(\delta)\) depends only on constant \(\delta\). This is (probably) because in this case we only need to constrain the sampled support, whereas the LMM lemma needs to constrain all \(n\) pure strategies.
\end{remark}
